\chapter{AI in Network Operations}\label{ch:ai}
\bp{5.1,5.2}

\begin{outcomes}
By the end of this chapter you will be able to:
\begin{tight}
  \item \textbf{Distinguish} generative AI from agentic AI, and describe what
        makes a system agentic.
  \item \textbf{Describe} where AI is genuinely used in network operations, and
        where it should not be trusted.
  \item \textbf{Select} an effective prompt, using the four components the
        blueprint names: data classification, output format, persona and
        instructions.
  \item \textbf{Judge} what data may safely be sent to an AI system, and sanitise
        what may not.
\end{tight}
\end{outcomes}

\begin{prereq}
\chref{management} for the five management approaches --- agentic systems sit on
top of them, not instead of them. \chref{ansible} for automation, which is what an
agent actually uses to act. \chref{aaa} for why credentials must not leave the
organisation.
\end{prereq}

\begin{keyterms}
generative AI $\bullet$ agentic AI $\bullet$ autonomy $\bullet$ tool use
$\bullet$ human in the loop $\bullet$ prompt $\bullet$ persona $\bullet$ data
classification $\bullet$ sanitisation $\bullet$ hallucination $\bullet$ blast
radius
\end{keyterms}

\begin{examnote}[Two topics, and the second one is unusual]
\tp{5.1} is \emph{describe the role of agentic AI in network operations}.

\tp{5.2} is \emph{select a prompt to send to a generative AI system to support
network operations considering prompt components such as data classification,
output format, persona, and instructions}.

Read \tp{5.2} carefully, because it tells you the question format. \emph{Select}
means you will be shown several candidate prompts and asked which is best ---
not asked to write one. And the four components are \textbf{named}, which means
the marking scheme is those four. A prompt that has all four beats one that
does not.

Note the order Cisco lists them in: \textbf{data classification comes first}.
That is not accidental, and Section~\ref{sec:dataclass} explains why.
\end{examnote}

\section{Generative and agentic}

\begin{definitionbox}[Generative AI]
A system that produces content --- text, code, configuration --- in response to a
prompt. It answers and stops. Every action taken as a result is taken by you.
\end{definitionbox}

\begin{definitionbox}[Agentic AI]
A system that pursues a \emph{goal} rather than answering a question. It
perceives the current state, plans, uses tools to act, observes the result, and
iterates --- choosing the intermediate steps itself.
\end{definitionbox}

\ccnafigh{%
\begin{tikzpicture}[font=\scriptsize]
  \tikzset{
    bx/.style={rounded corners=3pt, draw=#1, fill=#1!8, text=#1!45!black,
               line width=0.9pt, text width=2.5cm, align=center,
               minimum height=0.95cm, font=\tiny\bfseries},
    a/.style={-{Stealth[length=2.2mm]}, line width=0.9pt, neutral!70},
    ttl/.style={font=\scriptsize\bfseries, text=neutral!45!black, anchor=west}}

  \node[ttl] at (-0.3,1.55) {Generative};
  \node[bx=dom1] (g1) at (1.1,0.4)  {you ask};
  \node[bx=dom1] (g2) at (4.3,0.4)  {it answers};
  \node[bx=dom1] (g3) at (7.5,0.4)  {\textbf{you} act};
  \draw[a] (g1) -- (g2);  \draw[a] (g2) -- (g3);

  \node[ttl] at (-0.3,-1.15) {Agentic};
  \node[bx=dom5] (a1) at (1.1,-2.3) {goal set};
  \node[bx=dom5] (a2) at (4.3,-2.3) {observe\\state};
  \node[bx=dom5] (a3) at (7.5,-2.3) {plan};
  \node[bx=dom5] (a4) at (10.7,-2.3){\textbf{it} acts\\via tools};
  \draw[a] (a1) -- (a2);  \draw[a] (a2) -- (a3);  \draw[a] (a3) -- (a4);
  \draw[a, dom5!70] (a4) to[bend left=32]
      node[font=\tiny, text=dom5!45!black, midway, above=1pt]
      {check the result, iterate} (a2);
\end{tikzpicture}}%
{The difference is the loop and who closes it. A generative system hands you an
answer; an agentic one acts and checks its own work.}{agentic}

\begin{longtable}{@{}L{3.1cm}L{5.7cm}L{6.1cm}@{}}
\toprule
\rowcolor{primary!15}
\textbf{} & \textbf{Generative} & \textbf{Agentic} \\
\midrule
\endhead
Given & A prompt & A goal \\
\rowcolor{lightbg}
Produces & An answer & A sequence of actions \\
Tools & None. It only writes &
Calls APIs, runs show commands, opens tickets, applies configuration \\
\rowcolor{lightbg}
Steps & One & As many as it decides it needs \\
Who acts & \textbf{You} & \textbf{It} --- within whatever bounds you set \\
\rowcolor{lightbg}
Risk if wrong & A bad suggestion you can reject &
A change already made to a live network \\
\bottomrule
\end{longtable}

\begin{alertbox}[The one-sentence answer for the exam]
\textbf{Generative AI produces content; agentic AI takes action toward a goal,
using tools, and evaluates its own results.} Autonomy and tool use are the two
words that separate them.
\end{alertbox}

\section{Where it is actually used}

\begin{longtable}{@{}L{4.3cm}L{5.2cm}L{5.5cm}@{}}
\toprule
\rowcolor{primary!15}
\textbf{Role} & \textbf{What it does} & \textbf{Why it is suited to this} \\
\midrule
\endhead
Anomaly detection &
Learns normal behaviour and flags departures from it &
The volume of telemetry is far past what anyone can watch \\
\rowcolor{lightbg}
Event correlation &
Collapses a thousand alarms from one fault into one incident &
Pattern matching across sources, at speed \\
Root-cause suggestion &
Proposes the likely cause and the evidence for it &
It has seen the whole log, and you have seen the last screen \\
\rowcolor{lightbg}
Configuration drafting &
Produces candidate configuration and reviews existing configuration &
Syntax and boilerplate are exactly what it is good at \\
Documentation and summary &
Turns a change record or a packet capture into readable prose &
Tedious, low-risk, easily checked \\
\rowcolor{lightbg}
Capacity forecasting &
Projects growth from historical counters &
A statistical problem with plenty of data \\
Guided remediation &
Proposes a fix; with more autonomy, applies it &
This is the agentic end, and where care is needed \\
\bottomrule
\end{longtable}

\subsection{The autonomy ladder}

The useful question about an agentic system is not ``is it intelligent'' but
\textbf{how much is it allowed to do without asking}.

\begin{tight}
  \item \textbf{Recommend.} It tells you what it would do. You do it. No risk
        beyond a wasted suggestion.
  \item \textbf{Approve then act.} It prepares the change and waits for a human
        to approve. This is where most production deployments sit, and sensibly.
  \item \textbf{Act then report.} It makes the change and tells you afterwards.
        Appropriate for narrow, reversible, well-understood actions --- clearing
        an err-disabled port, restarting a stuck process.
  \item \textbf{Fully autonomous.} It acts within a defined scope with no
        notification. Appropriate for very little in a production network.
\end{tight}

\begin{conceptbox}[Blast radius is the design question]
Before granting autonomy, ask what the worst outcome is if the system is
confidently wrong. Clearing one err-disabled port has a blast radius of one port.
Changing a route map on a core router has a blast radius of the organisation.

The same technology can sit at different rungs for different actions, and a
well-designed deployment does exactly that. ``Human in the loop'' is not a
blanket policy; it is a decision made per action, based on reversibility and
reach.
\end{conceptbox}

\begin{alertbox}[What these systems are bad at, and you must know]
\begin{tight}
  \item \textbf{Being confidently wrong.} A fluent, plausible, incorrect answer
        is the characteristic failure. It does not signal uncertainty the way a
        colleague would.
  \item \textbf{Your network specifically.} Training data is the internet, not
        your topology. It does not know that VLAN~40 is the guest network here.
  \item \textbf{Currency.} It may recommend commands deprecated years ago, or
        miss ones added last year.
  \item \textbf{Accountability.} When a change breaks a site, a person is
        answerable. Change control, approval and rollback do not become optional
        because a machine proposed the change.
\end{tight}
\end{alertbox}

\section{Selecting a prompt}

This is topic \tp{5.2}, and it is a practical skill with four named parts.

\begin{longtable}{@{}L{3.4cm}L{5.4cm}L{6.2cm}@{}}
\toprule
\rowcolor{primary!15}
\textbf{Component} & \textbf{What it does} & \textbf{Weak against strong} \\
\midrule
\endhead
\textbf{Data classification} &
Decides what you may include, and what must be removed or replaced first &
Pasting a full running configuration with keys and public addresses, against
pasting the four relevant lines with addresses re-written \\
\rowcolor{lightbg}
\textbf{Persona} &
Sets the expertise and viewpoint to answer from &
``Explain VLANs'' against ``You are a senior network engineer working on Cisco
IOS XE 17.12'' \\
\textbf{Instructions} &
The actual task, with its constraints and its scope &
``Fix this'' against ``Identify why VLAN 30 clients get no address; do not
suggest changes to VLANs 10 or 20'' \\
\rowcolor{lightbg}
\textbf{Output format} &
The shape the answer must take, so it is usable &
``Tell me about it'' against ``Reply as a table of Cause, Evidence, Command to
confirm. No more than five rows'' \\
\bottomrule
\end{longtable}

\begin{worked}[The same question, asked twice]
\textbf{Weak:}

\begin{quote}
\RaggedRight\itshape Why is my DHCP not working? Here is my config:
[entire running configuration pasted, including \cmd{snmp-server community},
\cmd{username ... secret} and the site's public addressing]
\end{quote}

No persona, no output format, no scope --- and it has just sent credentials and
the organisation's real addressing to a third party. Even a good answer would not
justify that.

\textbf{Strong:}

\begin{quote}
\RaggedRight\itshape You are a senior network engineer working on Cisco IOS XE
17.12.

Clients in VLAN 30 receive APIPA addresses. VLANs 10 and 20 on the same
distribution switch lease correctly. The DHCP server is central and reached
across a routed core. Relevant configuration, with addresses substituted:

\medskip
\cmd{interface Vlan30 / ip address 10.30.0.1 255.255.255.0}

\medskip
List the three most likely causes, most likely first. For each give the single
show command that would confirm or eliminate it. Reply as a table with columns
Cause, Confirming command, Expected output if this is the fault. Do not propose
configuration changes.
\end{quote}

All four components are present, and each earns its place: the persona sets the
platform, the classification decided that addresses would be substituted and
credentials omitted, the instructions bound the scope, and the format makes the
answer directly actionable.
\end{worked}

\begin{examnote}[How to answer a ``select the best prompt'' question]
Score each candidate against the four components and pick the one with most.
Then apply three tie-breakers:

\begin{tight}
  \item \textbf{Eliminate anything containing credentials, keys or real public
        addressing} immediately, however good the rest of it is. Data
        classification is listed first for a reason and a prompt that leaks is
        wrong regardless of its other qualities.
  \item \textbf{Prefer specific over polite.} ``Please help me with my network''
        scores nothing. Length is not the measure --- \emph{specificity} is.
  \item \textbf{Prefer a stated output format.} It is the component most often
        missing from the distractors, so it frequently decides the answer.
\end{tight}
\end{examnote}

\subsection{Data classification, properly}\label{sec:dataclass}

\begin{alertbox}[Assume anything you send leaves your control]
A prompt sent to an external service has been disclosed to a third party. It may
be logged, retained, reviewed by humans, or used for training, depending on the
contract --- and deleting it later does not undo the disclosure.

\textbf{Never send:} passwords, hashes, SNMP community strings, pre-shared keys,
certificates, API tokens; personal data; real public IP addressing; anything
your organisation classifies as confidential or above.

\textbf{Sanitise first:} substitute addresses consistently (documentation ranges
\cmd{192.0.2.0/24}, \cmd{198.51.100.0/24}, \cmd{203.0.113.0/24} exist for exactly
this), replace hostnames, strip credential lines entirely, and include only the
configuration relevant to the question.

\textbf{And check the policy.} Many organisations permit an approved internal or
contracted tool and prohibit public ones. Knowing which applies to you is part
of the job now.
\end{alertbox}

\begin{pitfall}
\begin{tight}
  \item \textbf{Pasting a whole running configuration.} Slower, worse answers,
        and a disclosure.
  \item \textbf{Confusing generative with agentic.} Tool use and autonomy are the
        difference.
  \item \textbf{Omitting the output format.} You get an essay where you needed
        three commands.
  \item \textbf{No persona or platform.} You get generic advice, or IOS syntax
        for a platform you are not running.
  \item \textbf{Accepting a confident answer without verifying.} Check the
        command exists on your platform before typing it.
  \item \textbf{Granting autonomy without considering blast radius.}
  \item \textbf{Assuming it knows your network.} It knows the internet.
  \item \textbf{Treating an AI-proposed change as exempt from change control.}
\end{tight}
\end{pitfall}

\begin{breakfix}[The suggested command was accepted and took the site off the air.]
An engineer asks an assistant why a trunk is not carrying VLAN~50 and is told to
run \cmd{switchport trunk allowed vlan 50} on the uplink. They apply it. Every
VLAN except 50 immediately stops working across that link.

\begin{quote}
\RaggedRight\itshape \textbf{The prompt that was used:} ``VLAN 50 isn't going over
my trunk, what command fixes it?''
\end{quote}

\textbf{Diagnosis --- of the answer.} The command is real, correctly spelled and
valid. It also \emph{replaces} the allowed list rather than adding to it, so a
trunk carrying 10, 20 and 30 now carries only 50. The correct command was
\cmd{switchport trunk allowed vlan add 50} (\chref{trunking}).

\textbf{Diagnosis --- of the prompt}, which is the examinable part. Measure it
against the four components:

\begin{tight}
  \item \textbf{Persona} --- absent. Nothing established the platform or the
        expected level of care.
  \item \textbf{Instructions} --- absent. No mention that the trunk already
        carried other VLANs, which is the fact that makes the difference between
        the two commands. The model had no way to know it.
  \item \textbf{Output format} --- absent. A request for ``the command, plus what
        it changes, plus how to verify before and after'' would have surfaced the
        replace-versus-add behaviour in the answer itself.
  \item \textbf{Data classification} --- the one component that was, by accident,
        satisfied. Nothing sensitive was sent, because nothing was sent.
\end{tight}

The prompt asked a question whose correct answer depends on context it did not
supply. The response was reasonable given what it was told.

\textbf{The better prompt.}

\begin{quote}
\RaggedRight\itshape You are a senior network engineer working on Cisco IOS XE.
An 802.1Q trunk between two switches currently carries VLANs 10, 20 and 30. I
need to add VLAN 50 \textbf{without interrupting the existing three}. Give the
exact command, state what it changes, and give the show command to verify the
allowed list before and after. Warn me about any similar command that would be
destructive here.
\end{quote}

\textbf{The wider lesson, and it is the point of this chapter.} The fault was not
the tool's. A prompt missing three of its four components produced an answer that
was correct in general and wrong here, and it was applied to a production trunk
without verification. The habit that prevents this is not distrust of AI --- it
is the same habit that should apply to a command copied from a forum: understand
what it changes before you press enter, and have a way to check.
\end{breakfix}

\begin{lab}{Prompt engineering as an examinable skill}{Any generative AI assistant, plus a lab device}
\textbf{Platform note.} This lab needs no particular tool. Use whatever assistant
your institution permits --- and \textbf{check that policy first}, which is task
1 and is not a formality.

\begin{tasks}
  \item Find and read your institution's or employer's policy on sending data to
        AI services. Write down, in one sentence, what you are permitted to send.
  \item Take a real configuration from one of your lab devices. Produce a
        sanitised version fit to include in a prompt: substitute addressing into
        documentation ranges, rename hosts, remove every credential. List
        everything you removed and why.
  \item Write the weakest useful prompt you can for a real problem from
        \chref{tshootl2}. Record the answer.
  \item Rewrite it with all four components. Record the answer. Compare the two
        for specificity, correctness and usability.
  \item \textbf{Verify both answers on the device.} Note anything suggested that
        was wrong, deprecated, or not supported on your platform. This step is
        the one that matters.
  \item Ask a question whose correct answer depends on context you deliberately
        withhold --- as in this chapter's Break \& Fix. Confirm you can produce a
        confidently wrong answer, and identify which missing component caused it.
  \item Write four candidate prompts for one task, one of which is clearly best,
        and swap with a colleague. Score each other's against the four
        components. This is the exam question format.
  \item \textbf{Autonomy.} For each of these, place it on the four-rung ladder
        and justify: clearing an err-disabled port; adding a VLAN to an access
        switch; changing an OSPF cost on a core link; opening a ticket.
  \item \textbf{Extension.} Write the one-page guidance your team should follow
        for using AI assistants on network work --- what may be sent, what must
        be verified, and which actions may ever be automatic.
\end{tasks}
\end{lab}

\begin{checkpoint}
\begin{tightnum}
  \item State the difference between generative and agentic AI in one sentence.
  \item Name the four prompt components the blueprint lists, and say which is
        listed first and why.
  \item Give three categories of data that must never be sent to an external AI
        service.
  \item What is blast radius, and how does it govern how much autonomy an agent
        should have?
  \item An assistant gives a fluent, confident, wrong answer. What is this called
        and what is the defence?
  \item Name three genuine uses of AI in network operations.
\end{tightnum}
\end{checkpoint}

\begin{chaptersummary}
\begin{tight}
  \item Generative AI produces content in response to a prompt. \textbf{Agentic
        AI pursues a goal}, uses tools, acts, and evaluates its own results.
        Autonomy and tool use are the difference.
  \item Real roles: anomaly detection, event correlation, root-cause suggestion,
        configuration drafting, documentation, capacity forecasting, guided
        remediation.
  \item Autonomy is a ladder --- recommend, approve then act, act then report,
        fully autonomous --- chosen per action according to blast radius.
  \item The four prompt components: \textbf{data classification, output format,
        persona, instructions}. A prompt with all four beats one without.
  \item Data classification is listed first. Never send credentials, keys,
        personal data or real public addressing; sanitise into documentation
        ranges first; check your organisation's policy.
  \item In a ``select the prompt'' question: eliminate anything that leaks, prefer
        specificity over politeness, and look for the stated output format.
  \item These systems are confidently wrong, do not know your network, and may be
        out of date. Verify before you apply, and keep change control.
\end{tight}
\end{chaptersummary}

\begin{reviewq}
  \item \textbf{(MCQ)} What most distinguishes agentic AI from generative AI?
        \begin{tight}
          \item[(a)] It produces longer answers
          \item[(b)] It takes actions using tools and evaluates the results
          \item[(c)] It runs on the network device itself
          \item[(d)] It requires no prompt
        \end{tight}
  \item \textbf{(MCQ)} Which prompt component decides what you are permitted to
        include in the first place?
        \begin{tight}
          \item[(a)] Persona \quad (b)~Output format
          \item[(c)] Data classification \quad (d)~Instructions
        \end{tight}
  \item \textbf{(MCQ)} Which of these should never appear in a prompt sent to an
        external service?
        \begin{tight}
          \item[(a)] A VLAN number \quad (b)~An SNMP community string
          \item[(c)] An interface name \quad (d)~A documentation-range address
        \end{tight}
  \item \textbf{(Short)} Explain blast radius and use it to justify why clearing
        an err-disabled port may be automated while changing a core route map
        should not.
  \item \textbf{(Short)} An assistant recommends a command that is valid but
        destructive in your situation. Name the two prompt components whose
        absence caused this, and rewrite the prompt.
  \item \textbf{(Applied)} Write a prompt asking for help diagnosing an OSPF
        adjacency stuck in EXSTART, containing all four components, using only
        data that is safe to send. Then state how you would verify the answer
        before acting on it.
  \item \textbf{(Applied)} Your manager proposes letting an agentic system apply
        ACL changes automatically out of hours. Give your recommendation with
        three specific conditions you would require, referring to the autonomy
        ladder.
\end{reviewq}
